\section*{Capítulo \ref{ch:b2:wittig} --- Wittig y otras reacciones de formación de C=C}

\begin{solution}{exo:b2:wittig:1}
\ce{Ph3P + CH3CH2Br -> [Ph3PCH2CH3]+ + Br-} (sustitución bimolecular, calentando en un disolvente), y luego
butil-litio en éter anhidro o THF bajo gas inerte: \ce{[Ph3PCH2CH3]+ + BuLi -> Ph3P=CHCH3 + BuH + Li+}.
\end{solution}

\begin{solution}{exo:b2:wittig:2}
El iluro \ce{Ph3P=CH2} y el benzaldehído dan fenileteno (estireno), \ce{C6H5CH=CH2}, y
óxido de trifenilfosfina.
\end{solution}

\begin{solution}{exo:b2:wittig:3}
\ce{Ph3P=CHCOOEt} y \ce{Ph3P=CHCOCH3}: estabilizados. \ce{Ph3P=CHCH3}: no estabilizado, el que
necesita butil-litio (o hidruro de sodio, amiduro de sodio). \ce{Ph3P=CHC6H5}: estabilizado por un arilo, se prepara en la práctica con un
alcóxido o un hidróxido, y da mezclas de geometrías.
\end{solution}

\begin{solution}{exo:b2:wittig:4}
(\textit{E})-Pent-2-enoato de etilo, \ce{CH3CH2CH=CHCOOEt}, con fosfato de dietilo y sodio.
\end{solution}

\begin{solution}{exo:b2:wittig:5}
Wittig: el iluro \ce{Ph3P=CH2}, preparado a partir del yoduro de metiltrifenilfosfonio y butil-litio, da con la ciclohexanona solo
metilenciclohexano, con el doble enlace fuera del anillo. La deshidratación del 1-metilciclohexan-1-ol pasa
por un carbocatión terciario y pierde el hidrógeno que da el alqueno más sustituido (Zaitsev):
principalmente 1-metilciclohexeno, el isómero equivocado.
\end{solution}

\begin{solution}{exo:b2:wittig:6}
Ambos cortes son el mismo: benzaldehído y el iluro de bencilideno \ce{Ph3P=CHC6H5}. Ese iluro solo está
estabilizado por un arilo y da una mezcla (\textit{E})/(\textit{Z}). Para el (\textit{E}) puro, úsese la reacción
HWE del bencilfosfonato de dietilo con el benzaldehído.
\end{solution}

\begin{solution}{exo:b2:wittig:7}
Con \ce{Ph3P=CHCH3} se obtiene \ce{CH3CH2CH=CHCH3}, el (\textit{Z})-pent-2-eno, principalmente (\textit{Z}).
Con \ce{Ph3P=CHCOOEt}, el (\textit{E})-pent-2-enoato de etilo, principalmente (\textit{E}).
\end{solution}

\begin{solution}{exo:b2:wittig:8}
\ce{Ph3P=CH2 + C6H10O -> C7H12 + Ph3PO}. Masas: metilenciclohexano \qty{96.173}{g/mol},
óxido de trifenilfosfina \qty{278.291}{g/mol}; economía atómica $96.173/(96.173 + 278.291) =
96.173/374.464 \approx 25.7~\%$. Tres cuartas partes de la masa de los reactivos se van como óxido de fosfina.
% ledger: aw:C, aw:H, aw:O, aw:P
\end{solution}

\begin{solution}{exo:b2:wittig:9}
El fósforo, nucleófilo, desplaza el bromuro:
\begin{equation*}
\ce{P(OEt)3 + BrCH2COOEt -> [(EtO)3PCH2COOEt]+ + Br-}.
\end{equation*}
Después, el bromuro
ataca un carbono de un grupo etilo del ion fosfonio (sustitución bimolecular), liberando
bromoetano y formando el enlace \ce{P=O}:
\begin{equation*}
\ce{[(EtO)3PCH2COOEt]+ + Br- -> (EtO)2P(O)CH2COOEt + EtBr}.
\end{equation*}
\end{solution}

\begin{solution}{exo:b2:wittig:10}
\ce{OHC-CH=CH-CHO + 2Ph3P=CHC6H5 -> C6H5CH=CH-CH=CH-CH=CHC6H5 + 2Ph3PO}: dos equivalentes del
iluro de bencilideno (a partir del cloruro de benciltrifenilfosfonio y una base), uno por cada grupo aldehído.
Los nuevos dobles enlaces se forman como mezclas de geometrías; el isómero todo-(\textit{E}), el más
estable, se aísla por cristalización, o bien se isomeriza la mezcla.
\end{solution}

\begin{solution}{exo:b2:wittig:11}
\omterm{def:b2:enolates-aldol:aldol}{Aldólica}: benzaldehído con propanona en exceso e hidróxido de sodio (\cref{ch:b2:enolates-aldol});
reactivos baratos, agua como subproducto, producto (\textit{E}), pero la condensación puede producirse dos veces
(dibencilidenacetona) si la propanona no está en exceso. Wittig: benzaldehído con el \omterm{def:b2:wittig:ylide}{iluro estabilizado}
\ce{Ph3P=CHCOCH3}; (\textit{E}), un solo producto, sin doble condensación, pero el iluro es caro y
hay que eliminar el óxido de trifenilfosfina. Aquí la \omterm{def:b2:enolates-aldol:aldol}{aldólica} es la mejor vía.
\end{solution}

\begin{solution}{exo:b2:wittig:12}
Wittig: butanal con \ce{Ph3P=CHCH2CH2CH3} (a partir del 1-bromobutano y luego butil-litio), exento de sales:
principalmente (\textit{Z}); subproducto, el óxido de trifenilfosfina, y algo de isómero (\textit{E}). Alquino:
oct-4-ino (a partir del etino por dos alquilaciones con 1-bromopropano y amiduro de sodio) hidrogenado sobre un
catalizador de paladio envenenado (\cref{prop:b2:alkene-redox:syn-hydrogenation}): solo (\textit{Z}); los
subproductos son el bromuro de sodio y el amoníaco de las alquilaciones.
\end{solution}

\begin{solution}{pb:b2:wittig:1}
\textbf{1.} \ce{C23H46}, es decir, \ce{C_nH_{2n}} con $n = 23$: un doble enlace (o un anillo), aquí un
\ce{C=C}.
\textbf{2.} El doble enlace \ce{C9=C10}.
\textbf{3.} Nonanal \ce{CH3(CH2)7CHO} con el iluro de un 1-halotetradecano; o tetradecanal
\ce{CH3(CH2)12CHO} con el iluro de un 1-halononano.
\textbf{4.} 1-Bromotetradecano, \ce{CH3(CH2)13Br}.
\textbf{5.} La sustitución bimolecular por la voluminosa trifenilfosfina es rápida en un carbono primario no impedido
y no compite con la eliminación.
\textbf{6.} \ce{Ph3P + C14H29Br -> [Ph3PC14H29]+ + Br-}, bromuro de tetradeciltrifenilfosfonio.
\textbf{7.} Butil-litio (o hidruro de sodio, amiduro de sodio); el hidróxido es una base demasiado débil para un
iluro no estabilizado, y el agua que aporta protonaría el iluro.
\textbf{8.} No, solo lleva un grupo alquilo: principalmente (\textit{Z}) (\cref{prop:b2:wittig:stereo}).
\textbf{9.} Un anillo de cuatro eslabones \ce{P-C-C-O}: el fósforo unido al carbono que lleva la cadena
\ce{C13H27}, ese carbono al que lleva la cadena \ce{C8H17}, y este al oxígeno, unido
al fósforo.
\textbf{10.} El butil-litio y el iluro reaccionan con el agua y el oxígeno.
\textbf{11.} El doble enlace une el antiguo carbono carbonílico del nonanal (C9 del producto) al
antiguo carbono del iluro (C10), y ningún otro enlace se desplaza (\cref{prop:b2:wittig:regiospecific}).
\textbf{12.} Un carbocatión en C9, que puede perder un hidrógeno de C8 o de C10 y puede transponerse: una mezcla de
tricos-8-eno y tricos-9-eno, cada uno como (\textit{E}) y (\textit{Z}), sobre todo (\textit{E}).
\textbf{13.} \ce{Ph3P=CHC13H27 + C8H17CHO -> C8H17CH=CHC13H27 + Ph3PO}.
\textbf{14.} Objetivo \ce{C23H46}: \qty{322.621}{g/mol}; nonanal \ce{C9H18O}: \qty{142.242}{g/mol}; sal
\ce{C32H44BrP}: \qty{539.582}{g/mol}; \ce{Ph3PO}, \ce{C18H15OP}: \qty{278.291}{g/mol}.
\textbf{15.} El iluro \ce{C32H43P} pesa $539.582 - 80.912 = \qty{458.670}{g/mol}$; con el nonanal,
\qty{600.912}{g/mol} de reactivos para \qty{322.621}{g/mol} de producto: $322.621/600.912 \approx
53.7~\%$.
\textbf{16.} Por su polaridad: el hidrocarburo apolar atraviesa una columna corta de sílice con un
eluyente apolar mientras el óxido queda retenido; o bien el óxido cristaliza en un disolvente apolar frío.
\textbf{17.} El enlace \ce{P=O} formado, muy fuerte, pesa más que el enlace \ce{P-C} roto
(\cref{prop:b2:wittig:driving}).
\textbf{18.} Necesita un \omterm{def:b2:wittig:hwe}{fosfonato} con un grupo atractor de electrones en el carbono que reacciona, ausente
de una simple cadena alquílica, y da el (\textit{E})-alqueno.
\textbf{19.} Tricos-9-ino, \ce{CH3(CH2)7C#C(CH2)12CH3}, con hidrógeno sobre un catalizador de paladio
envenenado (paladio sobre carbonato de calcio con una sal de plomo y quinoleína).
\textbf{20.} Dec-1-ino, \ce{CH3(CH2)7C#CH}, desprotonado con amiduro de sodio, y luego 1-bromotridecano
\ce{CH3(CH2)12Br}: $8 + 2 + 13 = 23$ carbonos, con el triple enlace entre C9 y C10.
\textbf{21.} Una hidrogenación más daría tricosano, sin doble enlace; el catalizador envenenado
retiene el alquino mucho más fuertemente que el alqueno, que abandona la superficie antes de una segunda adición.
\textbf{22.} Ambas requieren dos etapas a partir de fragmentos disponibles. La vía del alquino da el isómero (\textit{Z})
limpiamente y solo deja sales; la vía de Wittig da algo de isómero (\textit{E}) y una gran masa de
óxido de trifenilfosfina.
\textbf{23.} 100~\%: \ce{C23H44 + H2 -> C23H46}, todos los átomos acaban en el producto.
\textbf{24.} $278.291/322.621 = \boldsymbol{\approx \qty{0.86}{g}}$ de óxido de trifenilfosfina por
gramo de feromona.
\end{solution}
